\section{Coupling constraints}\label{sec:constraints}

The corrected-grid bound of Proposition~\ref{prop:cellwise} rounds each coordinate
independently, and every rounded point is feasible because the domain is a
product. A coupling constraint, even one as sparse as $x_1=x_2$, destroys
this: independent rounding leaves the feasible set, and the minimum over
feasible grid points need not be a lower bound. This section shows that
totally unimodular (TU) coupling of the continuous variables, with data
aligned to a uniform mesh, admits a \emph{feasible} mean-preserving rounding.
Filtering, certificates, accuracy-independent state counts and exact output
then carry over. The price is visible in two places: the rounding is
correlated, so a bound on the full continuous Hessian replaces coordinate
curvature, and the mesh must be uniform, so the per-coordinate state count is
of order $\sqrt{n_c\kappa_c}$ instead of logarithmic in $n$. The last
subsection shows that both prices are intrinsic to this approach.

\subsection{Model}\label{sec:tu-model}

\begin{definition}[Aligned TU model]\label{def:tu-model}
Let $n_c\ge1$ and $n_d\ge0$. The feasible set is
\[
X=\bigl\{(x,z)\in\R^{n_c}\times\Z^{n_d}:\ \ell\le x\le u,\ z\in\mathcal Z,\
Ax+Bz\le b\bigr\},\qquad \mathcal Z=Z_1\times\dots\times Z_{n_d},
\]
with the following data.
\begin{enumerate}[label=(\roman*)]
\item $A\in\{0,\pm1\}^{m\times n_c}$ is totally unimodular; $B$ and $b$ are
rational.
\item $\ell<u$ componentwise (fixed continuous coordinates have been
substituted), and every $Z_k$ is a nonempty, explicitly listed finite set of
integers.
\item A rational \emph{mesh unit} $\eta>0$ is given such that
\begin{equation}\label{eq:tu-align}
\ell/\eta\in\Z^{n_c},\qquad u/\eta\in\Z^{n_c},\qquad
(b-Bz)/\eta\in\Z^m\quad\text{for every }z\in\mathcal Z .
\end{equation}
For integral $\ell,u,B,b$ one may take $\eta=1$. Write $D_0$ for the
denominator of $\eta$ in lowest terms.
\item $F(x,z)=\sum_{a\in\mathcal A}f_a(v_{S_a})$ with $v=(x,z)$, and a tree
decomposition with $N$ nodes and bags of size at most $p$ is given such that
every scope $S_a$ and every row support
$\{i:A_{ri}\ne0\}\cup\{k:B_{rk}\ne0\}$ is contained in some bag.
\item (Directional curvature.) A linear subspace $\mathcal K\subseteq\R^{n_c}$
and a rational $\bar L>0$ are given such that, for every $z\in\mathcal Z$,
$F(\cdot,z)$ is twice continuously differentiable on an open set containing
$[\ell,u]$ and
\begin{equation}\label{eq:tu-curv}
w^{\T}\nabla^2_{xx}F(x,z)\,w\le\bar L\norm w^2\qquad
(x\in[\ell,u],\ w\in\mathcal K).
\end{equation}
Either $\mathcal K=\R^{n_c}$, or $\mathcal K=\ker C$, where the rows of $C$
are the continuous parts $A_r$ of some \emph{equality pairs}, that is, pairs
of rows $r,r'$ with $(A_{r'},B_{r'},b_{r'})=-(A_r,B_r,b_r)$.
\end{enumerate}
\end{definition}

Equalities are thus written as two opposite inequalities; this preserves
total unimodularity. Only the continuous columns $A$ must be TU: the integer
columns $B$ are arbitrary, because the rounding below never changes $z$.
Throughout we write
\[
W=\max_i(u_i-\ell_i),\quad d=\max\{1,\max_k\lvert Z_k\rvert\},\quad
h_j=\eta\,2^{-j},\quad E_j=\frac{n_c\bar L h_j^2}{8},
\]
$S=\mathcal S=\argmin_XF$, $\OPT=\min_XF$, and $S_i=\{v_i:v\in S\}$ for the projection of
$S$ on coordinate $i$ (continuous or discrete). When a growth constant is
discussed, it is a set-growth constant
\begin{equation}\label{eq:tu-setgrowth}
F(v)-\OPT\ge g\dist(v,S)^2\qquad(v\in X),\qquad
\kappa_c=\max\{1,\bar L/g\},
\end{equation}
with the Euclidean distance on $\R^{n_c+n_d}$. Point growth at a unique
minimizer is the case $S=\{v^*\}$. Neither $g$ nor $S$ is used by any
algorithm below.

\begin{remark}[Directional versus coordinate curvature]\label{rem:tu-curv}
With $\mathcal K=\R^{n_c}$, \eqref{eq:tu-curv} implies the coordinate
curvature condition of Definition~\ref{def:curvature} with $L_i=\bar L$ for the
continuous coordinates, but not conversely. For a quadratic the largest
absolute row sum of $H_{xx}$, or any larger positive rational, is a valid
$\bar L$. For $\mathcal K=\ker C$ and a quadratic, $\bar L$ is valid if and
only if $V^{\T}(\bar LI-H_{xx})V\succeq0$ for a basis matrix $V$ of $\ker C$;
no orthonormality of $V$ is needed, since $w=Vc$ gives
$w^{\T}(\bar LI-H_{xx})w=c^{\T}V^{\T}(\bar LI-H_{xx})Vc$.
Alternatively, if $C$ has full row rank and
$\Pi=I-C^{\T}(CC^{\T})^{-1}C$, the largest absolute row sum of
$\Pi H_{xx}\Pi$ is valid, because $\Pi w=w$ for $w\in\ker C$ and the spectral
radius of a symmetric matrix is at most its largest absolute row sum. Both
checks are rational and change no factor scope. Example~\ref{ex:tu-fullcurv}
below shows that coordinate curvature alone is not enough.
\end{remark}

\begin{proposition}[Checking alignment without enumerating labels]\label{prop:tu-align}
Fix $z^0\in\mathcal Z$ and let $B_k$ be the $k$th column of $B$. Condition
\eqref{eq:tu-align} holds if and only if $\ell/\eta$ and $u/\eta$ are integral,
$(b-Bz^0)/\eta\in\Z^m$, and $B_k(t-z^0_k)/\eta\in\Z^m$ for every $k$ and
every $t\in Z_k$. This takes $1+\sum_k\lvert Z_k\rvert$ vector tests.
\end{proposition}

\begin{proof}
For every $z\in\mathcal Z$,
$(b-Bz)/\eta=(b-Bz^0)/\eta-\sum_kB_k(z_k-z^0_k)/\eta$, a sum of integral
vectors under the stated tests. Conversely, if \eqref{eq:tu-align} holds, the
tests are differences of two vectors $(b-Bz)/\eta$ for label vectors that
differ in at most one coordinate.
\end{proof}

\subsection{Feasible rounding on aligned cells}\label{sec:tu-rounding}

\begin{definition}[Level-$j$ domains]\label{def:tu-domain}
A \emph{level-$j$ domain} is a product
$D=D_1\times\dots\times D_{n_c}\times\Lambda_1\times\dots\times\Lambda_{n_d}$,
where each $D_i\subseteq[\ell_i,u_i]$ is a nonempty union of finitely many
closed intervals $[a,b]$, $a\le b$, with $a,b\in h_j\Z$, and each
$\Lambda_k\subseteq Z_k$ is nonempty. Its grid in coordinate $i$ is
$N_j(D_i)=D_i\cap h_j\Z$, its \emph{cells} are the intervals
$[t,t+h_j]\subseteq D_i$ with $t\in h_j\Z$, and its grid is
$G(D)=\prod_iN_j(D_i)\times\prod_k\Lambda_k$.
\end{definition}

Since $\ell,u\in\eta\Z^{n_c}\subseteq h_j\Z^{n_c}$, the box
$[\ell,u]\times\mathcal Z$ is a level-$0$ domain, and a level-$j$ domain is
also a level-$(j+1)$ domain. Components of $D_i$ may be single points.

\begin{lemma}[Feasible corner rounding]\label{lem:tu-round}
Let $D$ be a level-$j$ domain and $(x,z)\in X\cap D$. There is a random vector
$Y\in\R^{n_c}$ with finitely many values such that, almost surely,
$(Y,z)\in X\cap G(D)$ and $Y-x\in\mathcal K$, and
\begin{enumerate}[label=(\alph*)]
\item $\E Y=x$;
\item if $x_i\in h_j\Z$ then $Y_i=x_i$; otherwise, with
$t_i=h_j\lfloor x_i/h_j\rfloor$, the cell $[t_i,t_i+h_j]$ lies in $D_i$ and
$Y_i\in\{t_i,t_i+h_j\}$;
\item $\E\norm{Y-x}^2\le n_ch_j^2/4$.
\end{enumerate}
\end{lemma}

\begin{proof}
Write $h=h_j$ and $t_i=h\lfloor x_i/h\rfloor$, so $x=t+h\theta$ with
$\theta\in[0,1)^{n_c}$. Let $N_0=\{i:\theta_i=0\}=\{i:x_i\in h\Z\}$. If
$i\notin N_0$, then $x_i$ lies in a component $[a,b]$ of $D_i$ with
$a,b\in h\Z$ and $a<x_i<b$, so $a\le t_i$ and $t_i+h\le b$; the cell
$[t_i,t_i+h]$ lies in $D_i$.

Consider the polytope
\[
Q=\bigl\{\theta'\in\R^{n_c}:\ A\theta'\le\rho,\ 0\le\theta'\le\mathbf 1,\
\theta'_i\le0\ (i\in N_0)\bigr\},\qquad \rho=\frac{b-Bz-At}{h}.
\]
The right-hand side is integral: $(b-Bz)/h=2^j(b-Bz)/\eta\in\Z^m$ by
\eqref{eq:tu-align}, and $At/h\in\Z^m$ because $A$ is integral and
$t/h\in\Z^{n_c}$. Appending unit rows to a TU matrix gives a TU matrix
\cite[Section~19.1]{Schrijver1986}, so by the Hoffman--Kruskal theorem
\cite{HoffmanKruskal1956} every vertex of $Q$ is integral, hence lies in
$\{0,1\}^{n_c}$ and vanishes on $N_0$. Since $(x,z)\in X$, $\theta\in Q$, and the
polytope $Q$ is the convex hull of its vertices, so
$\theta=\sum_k\lambda_k\theta^k$ with vertices $\theta^k$ and convex weights
$\lambda_k$. Let $Y=t+h\theta^k$ with probability $\lambda_k$.

Then $\E Y=t+h\theta=x$. For each outcome,
$AY+Bz=At+hA\theta^k+Bz\le At+h\rho+Bz=b$. For $i\in N_0$, $Y_i=t_i=x_i$; for
$i\notin N_0$, $Y_i\in\{t_i,t_i+h\}$, which lies in the cell
$[t_i,t_i+h]\subseteq D_i\subseteq[\ell_i,u_i]$. Hence $(Y,z)\in X\cap G(D)$
and (b) holds. If $\mathcal K=\ker C$, each row $A_r$ of $C$ belongs to an
equality pair, so $A_rY+B_rz=b_r=A_rx+B_rz$ and $A_r(Y-x)=0$. Finally, for
$i\notin N_0$ the variable $Y_i$ takes two values with mean $x_i$, so
$\E(Y_i-x_i)^2=(x_i-t_i)(t_i+h-x_i)\le h^2/4$, and coordinates in $N_0$
contribute zero. This proves (c).
\end{proof}

The rounded coordinates are dependent: they follow a vertex decomposition of
the cell fiber, not a product law. This is why the next lemma needs
\eqref{eq:tu-curv} rather than coordinate curvature.

\begin{lemma}[Rounding allowance]\label{lem:tu-allow}
In the situation of Lemma~\ref{lem:tu-round},
\[
\E F(Y,z)\le F(x,z)+\frac{\bar L}2\E\norm{Y-x}^2\le F(x,z)+E_j .
\]
\end{lemma}

\begin{proof}
Fix an outcome and put $w=Y-x\in\mathcal K$. The segment from $x$ to $Y$ lies
in the box $[\ell,u]$, and $\phi(\tau)=F(x+\tau w,z)$ satisfies
$\phi''(\tau)=w^{\T}\nabla^2_{xx}F(x+\tau w,z)w\le\bar L\norm w^2$ by
\eqref{eq:tu-curv}. Taylor's formula with integral remainder gives
$\phi(1)\le\phi(0)+\phi'(0)+\bar L\norm w^2/2$, that is,
$F(Y,z)\le F(x,z)+\nabla_xF(x,z)^{\T}(Y-x)+\frac{\bar L}2\norm{Y-x}^2$.
Take expectations and use $\E(Y-x)=0$ and Lemma~\ref{lem:tu-round}(c).
\end{proof}

\begin{example}[Coordinate curvature is not enough]\label{ex:tu-fullcurv}
Let $0<h\le1/2$, $X=\{x\in[0,1]^2:x_1=x_2\}$ and
$F(x)=2x_1x_2-h(x_1+x_2)$. Both coordinate curvatures are zero, so the
corrections of Proposition~\ref{prop:cellwise} vanish. On $X$,
$F(s,s)=2s^2-2hs$ has minimum $\OPT=-h^2/2$ at $s=h/2$, whereas every feasible
point of the mesh-$h$ grid has value $F(kh,kh)=2h^2k(k-1)\ge0$. A
curvature-only correction would therefore certify the false bound $0$. The
Hessian has eigenvalues $\pm2$, so $\bar L=2$ is valid with
$\mathcal K=\R^2$, and $E=n_c\bar Lh^2/8=h^2/2$ gives the correct bound
$0-h^2/2=\OPT$; the allowance is attained.
\end{example}

\subsection{Filtering and certificates}\label{sec:tu-filter}

The tree computation of Lemma~\ref{lem:dp} applies to $F$ on a grid $G(D)$
after each row $r$ is assigned to a bag containing its support, as the table
$\iota_r(y)=0$ if $A_ry_x+B_ry_z\le b_r$ and $\iota_r(y)=+\infty$ otherwise.
The proof of Lemma~\ref{lem:dp} only adds and minimizes tables, so it holds
for values in $\R\cup\{+\infty\}$; when outgoing messages are formed by
subtracting one incoming term from a sum, one stores the sum of the finite
terms and the number of infinite ones. It returns
\[
v=\min\{F(y):y\in X\cap G(D)\},\qquad
M_i(t)=\min\{F(y):y\in X\cap G(D),\ y_i=t\},
\]
for every coordinate $i$ (continuous or discrete) and every grid value $t$,
and a minimizing grid point, with $\min\emptyset=+\infty$, using
$O\bigl(p(N+\lvert\mathcal A\rvert+m)K^p\bigr)$ table operations, where $K$
bounds the number of grid values per coordinate.

\paragraph{Algorithm TU-GRID$(\varepsilon)$.}
Set $D^{(0)}=[\ell,u]\times\mathcal Z$. For $j=0,1,2,\dots$:
\begin{enumerate}[label=(\arabic*)]
\item Compute $v_j$, a minimizer $y^{(j)}\in X\cap G(D^{(j)})$, and the
min-marginals $M^{(j)}_i$ on $G(D^{(j)})$. If $v_0=+\infty$, stop and report
$X=\emptyset$.
\item Put $U_j=v_j$, $\beta_j=v_j-E_j$ and $m_i(t)=M^{(j)}_i(t)-E_j$. If
$E_j\le\varepsilon$, return $y^{(j)}$ and $[\beta_j,U_j]$.
\item (Union filter.) For each continuous $i$, let $D^{(j+1)}_i$ be the union
of all cells $[t,t+h_j]$ of $D^{(j)}_i$ with
$\min\{m_i(t),m_i(t+h_j)\}\le U_j$ and of all single-point components
$\{t\}$ of $D^{(j)}_i$ with $m_i(t)\le U_j$. For each $k$, let
$\Lambda^{(j+1)}_k=\{t\in\Lambda^{(j)}_k:m_{k}(t)\le U_j\}$.
\end{enumerate}
The \emph{hull variant} replaces each $D^{(j+1)}_i$ by its convex hull. The
correction is the constant $E_j$, so $m_i$ and $\beta_j$ are the min-marginals
and the minimum of $Q=F-E_j$ on the feasible grid, in the notation of
Section~\ref{sec:grids}.

\begin{proposition}[Soundness]\label{prop:tu-sound}
For TU-GRID and for its hull variant, for every level $j$ that is reached:
\begin{enumerate}[label=(\alph*)]
\item if $X\ne\emptyset$ then $v_0<\infty$;
\item $\beta_j\le F(v)$ for every $v\in X\cap D^{(j)}$;
\item every $v\in(X\cap D^{(j)})\setminus D^{(j+1)}$ has $F(v)>U_j$;
\item $D^{(j+1)}$ is a level-$(j+1)$ domain that contains $S$ and $y^{(j)}$;
hence $y^{(j)}\in X\cap G(D^{(j+1)})$ and $U_{j+1}\le U_j$;
\item $\beta_j\le\OPT\le U_j\le\OPT+E_j$ and $U_j-\beta_j=E_j$.
\end{enumerate}
Moreover, a checker that recomputes every $v_j$ and $M^{(j)}_i$ from the
stored domains and verifies every filtering decision obtains
$\beta_J\le\OPT$ for the last level $J$ from (b) and (c) alone, without
using any minimizer, uniqueness or growth.
\end{proposition}

\begin{proof}
(a) Apply Lemma~\ref{lem:tu-round} with $j=0$ to any $(x,z)\in X$; its
outcomes are feasible points of $G(D^{(0)})$.

(b) Let $v=(x,z)\in X\cap D^{(j)}$ and let $Y$ be as in
Lemma~\ref{lem:tu-round}. Every outcome $(Y,z)$ is a feasible grid point, so
$v_j\le\E F(Y,z)\le F(x,z)+E_j$ by Lemma~\ref{lem:tu-allow}.

(c) Let $v=(x,z)\in(X\cap D^{(j)})\setminus D^{(j+1)}$ and let $Y$ be as in
Lemma~\ref{lem:tu-round}. If some label $z_k\notin\Lambda^{(j+1)}_k$, every
outcome has label $z_k$, so $M_k(z_k)\le F(Y,z)$ almost surely and
$m_k(z_k)\le\E F(Y,z)-E_j\le F(v)$; removal means $m_k(z_k)>U_j$. Otherwise
some $x_i\notin D^{(j+1)}_i$. If $x_i\in h_j\Z$, then $Y_i=x_i$ almost surely
and as before $m_i(x_i)\le F(v)$. Since $x_i\in D^{(j)}_i$, either $\{x_i\}$ is a
single-point component, which was removed, so $m_i(x_i)>U_j$; or $x_i$ is an
endpoint of at least one cell, all such cells were removed, and each removed
cell has both endpoint values above $U_j$, so again $m_i(x_i)>U_j$. If
$x_i\notin h_j\Z$, then $Y_i\in\{t_i,t_i+h_j\}$ for the cell
$[t_i,t_i+h_j]\subseteq D^{(j)}_i$ containing $x_i$, which was removed, and
$F(v)\ge\E F(Y,z)-E_j\ge\min\{m_i(t_i),m_i(t_i+h_j)\}>U_j$. In the hull variant
the retained set is larger, so (c) holds a fortiori.

(d) The retained sets are finite unions of closed intervals with endpoints in
$h_j\Z\subseteq h_{j+1}\Z$, as are their hulls. $X$ is compact and $F$ is
continuous, so $S\ne\emptyset$ when $X\ne\emptyset$. By induction
$S\subseteq D^{(j)}$; minimizers have value $\OPT\le F(y^{(j)})=U_j$, so by
(c) they are not removed. The point $y^{(j)}$ has value $U_j$, so it is not
removed either; its continuous coordinates lie in $h_j\Z\subseteq h_{j+1}\Z$,
hence $y^{(j)}\in X\cap G(D^{(j+1)})$ and $v_{j+1}\le v_j$. In particular
every $D^{(j+1)}_i$ and $\Lambda^{(j+1)}_k$ is nonempty.

(e) $\OPT\le U_j$ because $y^{(j)}\in X$. By (d) a minimizer lies in
$D^{(j)}$, so (b) gives $\beta_j\le\OPT$; rounding that minimizer as in (b)
gives $v_j\le\OPT+E_j$.

For the last claim, let $v\in X$. If $v\in D^{(J)}$, then $F(v)\ge\beta_J$ by
(b). Otherwise $v$ is removed at a first level $j<J$, and
$F(v)>U_j\ge U_J\ge\beta_J$ by (c); the inequality $U_j\ge U_J$ is checked
directly from the recomputed values. Hence $\OPT\ge\beta_J$.
\end{proof}

The certificate record consists of the data, a TU certificate for $A$
(a structural certificate such as a network or interval representation, or a
run of a polynomial-time TU recognition algorithm \cite[Chapter~20]{Schrijver1986}),
the curvature certificate for \eqref{eq:tu-curv}, $\eta$, the domains, the
tables, and the points $y^{(j)}$. As in Theorem~\ref{thm:certificate}, the
whole history is needed: the last grid bounds $F$ only on $X\cap D^{(J)}$.

\subsection{Accuracy-independent state counts}\label{sec:tu-states}

\begin{theorem}[State counts under set growth]\label{thm:tu-states}
Assume \eqref{eq:tu-setgrowth} and put $a_j=\sqrt{n_c\bar L/(4g)}\,h_j$. For
every level $j$ reached by TU-GRID:
\begin{enumerate}[label=(\alph*)]
\item every point of $D^{(j+1)}_i$ lies within $a_j+h_j$ of $S_i$, and every
label in $\Lambda^{(j+1)}_k$ lies within $a_j$ of $S_k$;
\item if a continuous projection $S_i$ has at most $r_i$ points, then
\[
\bigl\lvert N_{j+1}(D^{(j+1)}_i)\bigr\rvert\le
r_i\Bigl(5+\bigl\lfloor2\sqrt{n_c\bar L/g}\bigr\rfloor\Bigr)
\le r_i\bigl(5+\bigl\lfloor2\sqrt{n_c\kappa_c}\bigr\rfloor\bigr);
\]
\item for the hull variant the bound in (b) holds with $r_i=1$ when $S$ is a
single point.
\end{enumerate}
\end{theorem}

\begin{proof}
(a) A retained cell has an endpoint $t$ with $m_i(t)\le U_j$, a retained
single point $\{t\}$ has $m_i(t)\le U_j$, and a retained label $t$ has
$m_k(t)\le U_j$. In each case $M_i(t)<\infty$, so there is a feasible grid
point $w$ with $w_i=t$ and, by Proposition~\ref{prop:tu-sound}(e),
\[
F(w)=M_i(t)=m_i(t)+E_j\le U_j+E_j\le\OPT+2E_j .
\]
By \eqref{eq:tu-setgrowth}, $\dist(w,S)^2\le2E_j/g=a_j^2$, so some $s\in S$
has $\lvert t-s_i\rvert\le\norm{w-s}\le a_j$. Points of a cell are within $h_j$
of each of its endpoints.

(b) By (a), $D^{(j+1)}_i$ lies in the union over $\sigma\in S_i$ of the
intervals $[\sigma-a_j-h_j,\sigma+a_j+h_j]$. Each has length
$2a_j+2h_j=h_{j+1}(4a_j/h_j+4)$, so it contains at most
$\lfloor4a_j/h_j\rfloor+5$ points of $h_{j+1}\Z$, and
$4a_j/h_j=2\sqrt{n_c\bar L/g}$. Finally $\bar L/g\le\kappa_c$.

(c) In the hull variant the argument of (a) applies to the retained cells
before the hull is taken. If $S=\{v^*\}$, they lie in the interval
$[v^*_i-a_j-h_j,v^*_i+a_j+h_j]$, and hence so does their hull.
\end{proof}

The bound is per optimal coordinate value: different coordinates may be
near different minimizers, and no optimal combination is enumerated. There
can be as many as $\prod_ir_i\prod_k\lvert Z_k\rvert$ minimizers. The finite
projection hypothesis cannot be dropped for uniform grids: if $S_i$ contains a
nondegenerate interval $I$, then $I\subseteq D^{(j)}_i$ for every $j$ by
Proposition~\ref{prop:tu-sound}(d), so the level-$j$ grid has at least
$\lvert I\rvert/h_j-1$ points in coordinate $i$, which grows with the accuracy.

\begin{example}[Two minimizers per block: union versus hull]\label{ex:tu-union}
For one block take continuous $x,t,z\in[0,1]$ with the TU equality $x+t=z$
and
\[
F_b(x,t,z)=(2x-z)^2+\tfrac14z(1-z).
\]
Both terms are nonnegative on $X_b$, so $S_b=\{(0,0,0),(\tfrac12,\tfrac12,1)\}$
and $\OPT=0$. On $X_b$ put $w=x-z/2$, so $t=z/2-w$ and
$F_b=4w^2+z(1-z)/4$. If $z\le1/2$, the squared distance to $(0,0,0)$ is
$2w^2+\frac32z^2$ and $z(1-z)/4\ge z/8\ge z^2/4$, so
$F_b\ge4w^2+z^2/4\ge\frac16(2w^2+\frac32z^2)$. If $z\ge1/2$, the same
computation with $z'=1-z$ and the minimizer $(\frac12,\frac12,1)$ gives the
same inequality. Hence $g=1/6$. The Hessian in the order $(x,t,z)$ has rows
$(8,0,-4)$, $(0,0,0)$, $(-4,0,3/2)$, so $\bar L=12$ is valid. For $n_b$
independent blocks, $n_c=3n_b$, and the bags are the blocks ($p=3$). The
optimal set is the product of the block sets, so squared distances to it add
over blocks and $g=1/6$ and $\bar L=12$ are unchanged; every projection has
$r_i=2$ points, and there are $2^{n_b}$ minimizers. With $\eta=1$, the hull variant keeps $[0,1]$ in every $z$
coordinate at every level, because both $0$ and $1$ are optimal values, so
its level-$j$ grid has $2^j+1$ points there. The union variant has at most
$2(5+\lfloor2\sqrt{216n_b}\rfloor)$ points per coordinate at every level
$j\ge1$ by Theorem~\ref{thm:tu-states}.
\end{example}

\subsection{Complexity of certified approximation}\label{sec:tu-complexity}

\begin{theorem}[Certified approximation with TU coupling]\label{thm:tu-approx}
Assume Definition~\ref{def:tu-model} with factors that are explicit rational
quadratics or explicitly encoded rational polynomials of fixed degree, and
let $I$ be the binary input length, including $\eta$, $\bar L$ and the
decomposition. For $\varepsilon>0$, TU-GRID$(\varepsilon)$ either certifies
$X=\emptyset$ at level $0$, or stops at
\[
J=\min\{j\ge0:E_j\le\varepsilon\}\le
\max\Bigl\{0,\Bigl\lceil\tfrac12\log_2\frac{n_c\bar L\eta^2}{8\varepsilon}\Bigr\rceil\Bigr\}
\]
with a feasible rational point and an interval $[\beta_J,U_J]\ni\OPT$ of width
$E_J\le\varepsilon$, certified as in Proposition~\ref{prop:tu-sound}. All
numbers have bit length polynomial in $I+J$, and the number of table
operations is
\[
O\Bigl(p\,(N+\lvert\mathcal A\rvert+m)\sum_{j=0}^JK_j^p\Bigr),\qquad
K_0\le\max\{d,W/\eta+1\}.
\]
Under \eqref{eq:tu-setgrowth} with $\lvert S_i\rvert\le r$ for every continuous
$i$, $K_j\le K=\max\{d,r(5+\lfloor2\sqrt{n_c\kappa_c}\rfloor)\}$ for every
$j\ge1$. The algorithm uses neither $g$ nor $r$.
\end{theorem}

\begin{proof}
Correctness is Proposition~\ref{prop:tu-sound}, and the width is
$U_J-\beta_J=E_J$. The bound on $J$ follows from $E_j=n_c\bar L\eta^24^{-j}/8$.
The table count is the tree computation above, summed over levels; the
level-$0$ grid has $(u_i-\ell_i)/\eta+1\le W/\eta+1$ points in a continuous
coordinate and at most $d$ labels in a discrete one, and labels are never
added. The bound $K$ is Theorem~\ref{thm:tu-states}(b) together with
$\lvert\Lambda^{(j)}_k\rvert\le d$. Grid points at level $j$ lie in
$h_j\Z\cap[\ell,u]$, so their coordinates are rationals with denominators
dividing $D_02^j$ and numerators bounded by $D_02^j\max_i\max\{\lvert\ell_i\rvert,\lvert u_i\rvert\}$.
Explicit fixed-degree factors evaluated at such points have polynomial bit
length, and every table entry is a sum of at most $\lvert\mathcal A\rvert$ factor
values or a minimum of such sums.
\end{proof}

Without growth the same certificate is produced, but the tables can be as
large as $(W2^j/\eta+1)^p$; by minimality of $J$ this is
$O\bigl((W\sqrt{n_c\bar L/(2\varepsilon)}+1)^p\bigr)$ at the last level. With
growth, for fixed $p$ and polynomially bounded $\kappa_c$, $W/\eta$, $d$ and $r$, the
algorithm is polynomial. It is not fixed-parameter tractable in
$(p,\kappa_c)$: $K^p$ contains the factor $n_c^{p/2}$, and
Proposition~\ref{prop:tu-tight} below shows that this factor is really
incurred. Large integral capacities make the level-$0$ grid
pseudopolynomial; a larger common unit $\eta$, certified by
Proposition~\ref{prop:tu-align}, reduces that cost to $W/\eta+1$.

\subsection{Exact output for quadratic objectives}\label{sec:tu-exact}

In this subsection $F(v)=\frac12v^{\T}Hv+q^{\T}v+c_0$ is a quadratic with
rational data on the model of Definition~\ref{def:tu-model}. Let $\Delta$ be a
positive common denominator of the monomial coefficients $H_{ii}/2$,
$H_{ik}$ $(i<k)$, $q_i$ and $c_0$, and put $P=\Delta H$, which is integral and
symmetric. Write $P_{xx}$ and $P_{xz}$ for its continuous and mixed blocks and
\[
C_P=\max\Bigl\{1,\max_{i,k\le n_c}\lvert(P_{xx})_{ik}\rvert\Bigr\},\quad
R=(2n_cC_P)^{n_c},\quad \Omega=\Delta(D_0R)^2,\quad
\tau=\frac1{4m'D_0R},
\]
where $m'=m+2n_c$. These numbers have bit length polynomial in $I$ and are
computed from the original data only. If $z$ is integral and
$x\in\rho^{-1}\Z^{n_c}$ for a positive integer $\rho$, then
$\Delta\rho^2F(x,z)\in\Z$, since $\Delta$ clears every monomial coefficient.

For $z\in\mathcal Z$ let
\[
M=\begin{pmatrix}A\\ I\\ -I\end{pmatrix}\in\{0,\pm1\}^{m'\times n_c},\qquad
d_z=\begin{pmatrix}b-Bz\\ u\\ -\ell\end{pmatrix}\in D_0^{-1}\Z^{m'},\qquad
\mathcal P_z=\{x:Mx\le d_z\},
\]
so that $\{x:(x,z)\in X\}=\mathcal P_z$. Put $\phi_z(x)=F(x,z)$. Then
$\Delta\nabla\phi_z(x)=P_{xx}x+c_z$ with $c_z=P_{xz}z+\Delta q_x\in\Z^{n_c}$. For
a row set $J$ write $M_J$ for the corresponding rows.

\begin{definition}[Stationary face polytope]\label{def:tu-statpoly}
For $s=(s_x,z)\in S$ let $J_0=\{r:M_rs_x=(d_z)_r\}$ be its active rows and
\[
\mathcal P(s)=\bigl\{x\in\mathcal P_z:\ M_{J_0}x=(d_z)_{J_0},\ \
\nabla\phi_z(x)\perp\ker M_{J_0}\bigr\}.
\]
\end{definition}

\begin{lemma}[Stationary face polytopes]\label{lem:tu-statpoly}
Let $s=(s_x,z)\in S$ and $J_0$ as above.
\begin{enumerate}[label=(\alph*)]
\item $\nabla\phi_z(s_x)\perp\ker M_{J_0}$, and $w^{\T}P_{xx}w\ge0$ for
$w\in\ker M_{J_0}$.
\item $\mathcal P(s)$ is a nonempty polytope and $\mathcal P(s)\times\{z\}\subseteq S$.
\item Every vertex $v$ of $\mathcal P(s)$ lies in $(D_0\delta)^{-1}\Z^{n_c}$ for an
integer $1\le\delta\le R$.
\end{enumerate}
\end{lemma}

\begin{proof}
(a) For $w\in\ker M_{J_0}$ and small $\lvert\tau\rvert$, $s_x+\tau w\in\mathcal P_z$:
active rows stay active and inactive rows keep positive slack. The function
$\tau\mapsto\phi_z(s_x+\tau w)=\OPT+\tau\nabla\phi_z(s_x)^{\T}w+\frac{\tau^2}{2}w^{\T}H_{xx}w$
has a minimum at $\tau=0$, which gives both claims.

(b) By (a), $s_x\in\mathcal P(s)$. The condition
$\nabla\phi_z(x)\perp\ker M_{J_0}$ reads $N^{\T}(P_{xx}x+c_z)=0$ for a basis
matrix $N$ of $\ker M_{J_0}$, which is linear, and $\mathcal P(s)\subseteq\mathcal P_z$
is bounded. For $x\in\mathcal P(s)$ put $w=x-s_x\in\ker M_{J_0}$. Both
$\nabla\phi_z(x)$ and $\nabla\phi_z(s_x)$ are orthogonal to $\ker M_{J_0}$, hence so
is their difference $H_{xx}w$, and
\[
\phi_z(x)-\phi_z(s_x)=\nabla\phi_z(s_x)^{\T}w+\tfrac12w^{\T}H_{xx}w=0 .
\]
Since $(x,z)\in X$, it is a minimizer.

(c) Let $J_v\supseteq J_0$ be the rows active at $v$, let $E$ be a set of
$r$ linearly independent rows spanning the rows of $M_{J_v}$, and let $e$ be
the corresponding entries of $d_z$. First, $P_{xx}$ is positive definite on
$\ker E=\ker M_{J_v}\subseteq\ker M_{J_0}$. Indeed, let $w\in\ker E$ with
$w^{\T}P_{xx}w=0$. By (a) the quadratic form of $P_{xx}$ is nonnegative on
$\ker M_{J_0}$ and vanishes at $w$, so for $y\in\ker M_{J_0}$ and all real
$\tau$, $0\le(w+\tau y)^{\T}P_{xx}(w+\tau y)=2\tau y^{\T}P_{xx}w+\tau^2y^{\T}P_{xx}y$,
whence $P_{xx}w\perp\ker M_{J_0}$. Then $v\pm\tau w\in\mathcal P(s)$ for small
$\tau>0$: rows in $J_v$ stay active, the other rows keep positive slack, and
the stationarity condition is preserved. As $v$ is a vertex, $w=0$.

Next, $\nabla\phi_z(v)\perp\ker M_{J_0}\supseteq\ker E$, so
$P_{xx}v+c_z=E^{\T}\mu$ for some $\mu$. Thus $(v,-\mu)$ solves
\[
K\begin{pmatrix}v\\ -\mu\end{pmatrix}=\begin{pmatrix}-c_z\\ e\end{pmatrix},\qquad
K=\begin{pmatrix}P_{xx}&E^{\T}\\ E&0\end{pmatrix}.
\]
$K$ is nonsingular: if $P_{xx}w+E^{\T}\nu=0$ and $Ew=0$, then
$w^{\T}P_{xx}w=-(Ew)^{\T}\nu=0$, so $w=0$, and then $E^{\T}\nu=0$ gives $\nu=0$
because $E$ has full row rank. The matrix $K$ is integral and the right-hand
side lies in $D_0^{-1}\Z^{n_c+r}$, so Cramer's rule gives
$v\in(D_0\delta)^{-1}\Z^{n_c}$ with $\delta=\lvert\det K\rvert\ge1$. By Hadamard's
inequality for rows, $\delta$ is at most the product of the row norms of $K$.
A row of $(P_{xx}\ E^{\T})$ has $n_c$ entries of absolute value at most $C_P$
and $r\le n_c$ entries in $\{0,\pm1\}$, so its norm is at most
$\sqrt{2n_c}\,C_P$; a row of $(E\ 0)$ is a row of $M$ and has norm at most
$\sqrt{n_c}$. Hence
$\delta\le(2n_cC_P^2)^{n_c/2}n_c^{r/2}\le(\sqrt2\,n_cC_P)^{n_c}\le R$.
\end{proof}

\begin{corollary}[Rational height with TU coupling]\label{cor:tu-height}
Some global minimizer has continuous coordinates in $\rho^{-1}\Z$ for an
integer $1\le\rho\le D_0R$, and every isolated global minimizer has this
property. The optimal value is $\OPT=a/W$ in lowest terms with $1\le W\le\Omega$.
\end{corollary}

\begin{proof}
For $s\in S$ the polytope $\mathcal P(s)$ has a vertex, which is a minimizer by
Lemma~\ref{lem:tu-statpoly}(b) and has the stated height by (c). If $s$ is
isolated in $S$, the convex set $\mathcal P(s)\times\{z\}\subseteq S$ contains
no other point near $s$, so $\mathcal P(s)=\{s_x\}$ and $s_x$ is its vertex.
Evaluating $F$ at a minimizer of height $\rho\le D_0R$ gives
$\Delta\rho^2\OPT\in\Z$, so $W\le\Delta(D_0R)^2=\Omega$.
\end{proof}

Corollary~\ref{cor:tu-height} is an explicit-constant instance of the
classical fact that quadratic programs over polyhedra have optimal solutions
of polynomial size \cite{Vavasis1990}; for mixed-integer quadratic programs see
\cite{DelPiaDeyMolinaro2017}. The constant uses only the continuous Hessian
block and the $\{0,\pm1\}$ structure of $M$.

\begin{proposition}[Exact acceptance]\label{prop:tu-accept}
Let $\beta\le\OPT$ and $v\in X$ with $F(v)=a/W$ in lowest terms. If
$F(v)-\beta<1/(\Omega W)$, then $v\in S$.
\end{proposition}

\begin{proof}
If $F(v)>\OPT$, the two rationals have denominators $W$ and at most $\Omega$
(Corollary~\ref{cor:tu-height}), so $F(v)-\OPT\ge1/(\Omega W)$, contradicting
$\beta\le\OPT$.
\end{proof}

\paragraph{Procedure TU-REC$(y,z)$.}
Given $(y,z)\in X$, let $J=\{r:(d_z)_r-M_ry\le\tau\}$ and let $N$ be a basis
matrix of $\ker M_J$, computed by Gaussian elimination. Find, by exact linear
programming, a point of
\[
\Pi=\bigl\{x\in\mathcal P_z:\ M_Jx=(d_z)_J,\ N^{\T}(P_{xx}x+c_z)=0\bigr\},
\]
or report failure. All coefficients come from the original data, the index
set $J$ and the integer vector $z$, so this takes time polynomial in $I$
\cite{GrotschelLovaszSchrijver1988}; the large denominators of
$y$ enter only the comparisons that define $J$.

\begin{lemma}[Snapping recovery on TU faces]\label{lem:tu-snap}
If $(y,z)\in X$ and $\dist((y,z),S)\le\tau/(2\sqrt{n_c})$, then $\Pi\ne\emptyset$
and $\Pi\times\{z\}\subseteq S$.
\end{lemma}

\begin{proof}
Let $s$ be a point of $S$ nearest to $(y,z)$. Since $\norm{(y,z)-s}<1$, the
integer parts agree, so $s=(s_x,z)$ and $\norm{y-s_x}\le\tau/(2\sqrt{n_c})$. Let
$J_0$ be the rows active at $s_x$. Every row of $M$ has norm at most $\sqrt{n_c}$,
so for $r\in J_0$, $0\le(d_z)_r-M_ry=M_r(s_x-y)\le\tau/2$; hence $J_0\subseteq J$.
For $r\in J\setminus J_0$ the slack at $s_x$ is at most $\tau+\tau/2$. On
$\mathcal P(s)$ the linear function
\[
\psi(x)=\sum_{r\in J\setminus J_0}\bigl((d_z)_r-M_rx\bigr)
\]
is nonnegative and $\psi(s_x)\le m'\cdot\frac32\tau=\frac3{8D_0R}<\frac1{D_0R}$.
Its minimum over the polytope $\mathcal P(s)$ is attained at a vertex $v$. By
Lemma~\ref{lem:tu-statpoly}(c), $v\in(D_0\delta)^{-1}\Z^{n_c}$ with $\delta\le R$; as $M$
is integral and $d_z\in D_0^{-1}\Z^{m'}$, $\psi(v)\in(D_0\delta)^{-1}\Z$. Since
$0\le\psi(v)<1/(D_0R)\le1/(D_0\delta)$, $\psi(v)=0$.

So $v$ satisfies every row of $J\setminus J_0$ with equality, and every row of
$J_0$ because $v\in\mathcal P(s)$. Moreover $\ker M_J\subseteq\ker M_{J_0}$ and
$\nabla\phi_z(v)\perp\ker M_{J_0}$, so $N^{\T}(P_{xx}v+c_z)=0$. Hence $v\in\Pi$.
For any $x\in\Pi$ let $w=x-v\in\ker M_J$. Both $\nabla\phi_z(x)$ and
$\nabla\phi_z(v)$ are orthogonal to $\ker M_J$, hence so is $H_{xx}w$, and
$\phi_z(x)=\phi_z(v)+\nabla\phi_z(v)^{\T}w+\frac12w^{\T}H_{xx}w=\phi_z(v)=\OPT$,
the last equality by Lemma~\ref{lem:tu-statpoly}(b). Since $x\in\mathcal P_z$,
$(x,z)\in S$.
\end{proof}

No nonsingularity, strict complementarity or sign condition on multipliers is
used: the condition $N^{\T}(P_{xx}x+c_z)=0$ allows multipliers of either
sign. A point returned from a distant $y$ can be a nonglobal stationary point;
for $F(x,t)=x-x^2$ on $\{0\le x=t\le1\}$ and $y=(\frac12,\frac12)$, the procedure
returns $(\frac12,\frac12)$ with value $\frac14>\OPT=0$. Acceptance therefore always
goes through Proposition~\ref{prop:tu-accept}.

\paragraph{Algorithm TU-EXACT.}
Run TU-GRID without a stopping accuracy. At every level $j$, apply TU-REC to
$y^{(j)}$; if it returns $\hat x$, let $\hat v=(\hat x,z^{(j)})$, where
$z^{(j)}$ is the label vector of $y^{(j)}$, and return $\hat v$ if
$F(\hat v)-\beta_j<1/(\Omega W)$, where $W$ is the reduced denominator of
$F(\hat v)$.

\begin{theorem}[Exact output with TU coupling]\label{thm:tu-exact}
Let $F$ be a rational quadratic on a nonempty set $X$ as in
Definition~\ref{def:tu-model}.
\begin{enumerate}[label=(\alph*)]
\item Every point returned by TU-EXACT is a global minimizer, and
$F(\hat v)=\OPT$. This uses neither uniqueness nor growth.
\item TU-EXACT terminates.
\item Under \eqref{eq:tu-setgrowth} it terminates by the first level $j$ with
$E_j\le\min\{g\tau^2/(4n_c),\,1/(2\Omega^2)\}$, hence after at most
\[
J_{\rm ex}=\max\Bigl\{0,\Bigl\lceil\log_2\frac{\eta n_c\sqrt{\kappa_c}}{\tau}\Bigr\rceil,
\Bigl\lceil\log_2\bigl(\eta\Omega\sqrt{n_c\bar L}\bigr)\Bigr\rceil\Bigr\}
=\poly(I)+O(\log\kappa_c)
\]
levels. If moreover every continuous projection $S_i$ has at most $r$ points,
the table operations are as in Theorem~\ref{thm:tu-approx} with $J$ replaced
by $J_{\rm ex}$; the LP calls add a polynomial cost per level.
\end{enumerate}
\end{theorem}

\begin{proof}
(a) $\beta_j\le\OPT$ by Proposition~\ref{prop:tu-sound}(e), and $\hat v\in X$ is
checked, so Proposition~\ref{prop:tu-accept} applies; then $F(\hat v)=\OPT$.

(b) By Proposition~\ref{prop:tu-sound}(e), $F(y^{(j)})=U_j\to\OPT$. We claim
$\dist(y^{(j)},S)\to0$. Otherwise a subsequence stays at distance at least
some $\epsilon>0$ from $S$ and, $X$ being compact, converges to some
$\bar v\in X$ with $F(\bar v)=\OPT$ and $\dist(\bar v,S)\ge\epsilon$, which is
impossible. So for all large $j$, $\dist(y^{(j)},S)\le\tau/(2\sqrt{n_c})$ and
$E_j<1/\Omega^2$. Then Lemma~\ref{lem:tu-snap} shows that TU-REC returns
$\hat x$ with $\hat v\in S$; its value is $\OPT=a/W$ with $W\le\Omega$, and
$F(\hat v)-\beta_j=\OPT-\beta_j\le U_j-\beta_j=E_j<1/\Omega^2\le1/(\Omega W)$, so
$\hat v$ is returned.

(c) By \eqref{eq:tu-setgrowth} and Proposition~\ref{prop:tu-sound}(e),
$g\dist(y^{(j)},S)^2\le U_j-\OPT\le E_j$, so the two conditions used in (b)
hold as soon as $E_j\le g\tau^2/(4n_c)$ and $E_j\le1/(2\Omega^2)$. Since
$g\ge\bar L/\kappa_c$ and $E_j=n_c\bar Lh_j^2/8$, it suffices that
$h_j\le\tau/(n_c\sqrt{\kappa_c})$ and $h_j\le2/(\Omega\sqrt{n_c\bar L})$, which hold
for $j\ge J_{\rm ex}$ because $h_j=\eta\,2^{-j}$. The quantities $\log_2(1/\tau)$, $\log_2\Omega$,
$\log_2\eta$ and $\log_2\bar L$ are polynomial in $I$. The last statement
follows from Theorem~\ref{thm:tu-states}.
\end{proof}

\begin{remark}[Unique minimizers without linear programming]\label{rem:tu-cf}
If $S=\{v^*\}$, the hull variant can return $v^*$ by continued fractions
alone. By Corollary~\ref{cor:tu-height}, $v^*_x\in\rho^{-1}\Z^{n_c}$ with
$\rho\le D_0R$, and by Theorem~\ref{thm:tu-states}(a) the retained hull in
coordinate $i$ contains $v^*_i$ and has length at most $2a_j+2h_j$. Two
distinct rationals with denominators at most $D_0R$ differ by at least
$(D_0R)^{-2}$, so once $2a_j+2h_j<(D_0R)^{-2}$ and $a_j<1$, the rational of
least denominator in each continuous hull is $v^*_i$, and each label set is
$\{v^*_k\}$. Once also $E_j<\Omega^{-2}$, the candidate passes the test of
Proposition~\ref{prop:tu-accept}; this check, not the reconstruction, carries
the proof.
\end{remark}

\subsection{What prevents fixed-parameter tractability}\label{sec:tu-limits}

The box algorithm of Section~\ref{sec:growth} is fixed-parameter tractable in
$(p,\kappa)$ because graded grids cover a localization radius of order
$\sqrt{n\kappa}\,h$ with $O(\theta^{-1}\log n)$ points. With TU coupling the
radius is the same, $a_j+h_j$ in Theorem~\ref{thm:tu-states}, but the mesh is
uniform. The next proposition shows that the resulting $\Theta(\sqrt{n_c})$
states per coordinate are actually retained; the two after it show that
graded or otherwise non-aligned grids cannot simply be used instead, because
a correction computed from curvature and the grids alone is then not a valid
lower bound.

\begin{proposition}[Uniform meshes retain $\Theta(\sqrt{n_c})$ points]\label{prop:tu-tight}
Let $n_c\ge p\ge1$, $\bar L>0$, $\eta=1$, and
\[
X=\{x\in[0,1]^{n_c}:x_1+\dots+x_p\le p\},\qquad
F(x)=\frac{\bar L}2\Bigl\lVert x-\tfrac12\mathbf 1\Bigr\rVert^2 .
\]
The single row is TU, every admissible decomposition has a bag containing
$\{1,\dots,p\}$, $\mathcal K=\R^{n_c}$ with $\bar L$ is valid, and point
growth holds with $g=\bar L/2$, so $\kappa_c=2$. Let
$k=\lfloor\sqrt{n_c}/2\rfloor$. For every level $j\ge1$ with
$(k+1)h_j\le\frac12$, both variants of TU-GRID retain
\[
D^{(j+1)}_i=\bigl[\tfrac12-(k+1)h_j,\ \tfrac12+(k+1)h_j\bigr]\qquad(i=1,\dots,n_c),
\]
whose level-$(j+1)$ grid has $4k+5\ge2\sqrt{n_c}+1$ points. Hence the bag
containing $\{1,\dots,p\}$ has at least $(2\sqrt{n_c}+1)^p$ table entries at
every such level, while Theorem~\ref{thm:tu-states} gives the upper bound
$5+\lfloor2\sqrt{2n_c}\rfloor$ per coordinate.
\end{proposition}

\begin{proof}
The row is redundant on $[0,1]^{n_c}$, so $X$ is the box and $F$ has the
unique minimizer $\frac12\mathbf 1$, with $F(x)-\OPT=\frac{\bar L}2\norm{x-\frac12\mathbf 1}^2$.
Consider a level $j\ge1$ and a level-$j$ domain $D$ whose components are
intervals $D_i$ containing $\frac12$. As $\frac12\in h_j\Z$, the minimizer is a
grid point, so $v_j=U_j=0$. For $t\in N_j(D_i)$, setting the other
coordinates to $\frac12$ gives $M_i(t)=\frac{\bar L}2(t-\frac12)^2$, and no
grid point does better since $F(y)\ge\frac{\bar L}2(y_i-\frac12)^2$. Thus
$m_i(t)\le U_j$ if and only if $\lvert t-\frac12\rvert\le\frac{\sqrt{n_c}}2h_j$, that
is, $t=\frac12+ih_j$ with $\lvert i\rvert\le k$. If
$D_i\supseteq[\frac12-(k+1)h_j,\frac12+(k+1)h_j]$, the retained cells are exactly
those with an endpoint $\frac12+ih_j$, $\lvert i\rvert\le k$, and their union is the
displayed interval, which equals its hull. That interval contains
$[\frac12-(k+1)h_{j+1},\frac12+(k+1)h_{j+1}]$, so the hypothesis propagates.

It remains to start the induction. At level $0$ the grid is $\{0,1\}^{n_c}$,
$v_0=n_c\bar L/8=E_0$, and every $m_i(t)=0\le U_0$, so $D^{(1)}=[0,1]^{n_c}$. At
a level $j\ge1$ with $(k+1)h_j\ge\frac12$ and $D^{(j)}=[0,1]^{n_c}$, the computation
above keeps the cells with an endpoint within $kh_j$ of $\frac12$, whose union
is $[0,1]$. Hence $D^{(j)}=[0,1]^{n_c}$ up to the first level with
$(k+1)h_j\le\frac12$, which contains the required interval. The interval of
length $2(k+1)h_j$ has $4(k+1)+1$ points of $h_{j+1}\Z$, and
$4k+5\ge4(\sqrt{n_c}/2-1)+5$.
\end{proof}

The instance is a box problem in disguise, and the graded grids of
Section~\ref{sec:growth} solve it with $O(\log n)$ points per coordinate. The
proposition therefore isolates the uniform mesh, not the constraint, as the
source of the factor $n_c^{p/2}$. Uniformity, however, is forced by the
constraints, as follows.

\begin{proposition}[Curvature-only corrections need aligned grids]\label{prop:tu-misaligned}
Let $G_1,G_2\subseteq[0,1]$ be finite sets containing $0$ and $1$, let
$d_i:G_i\to[0,\infty)$ be arbitrary, let $L>0$, and let $m\in(0,1)$ satisfy
\begin{equation}\label{eq:tu-gap}
L(\nu-m)^2>d_1(\nu)+d_2(\nu)\qquad\text{for every }\nu\in G_1\cap G_2 .
\end{equation}
For $\lambda\ge0$ let $X=\{x\in[0,1]^2:x_1-x_2\le0\}$ and
$F_\lambda(x)=\frac L2\bigl((x_1-m)^2+(x_2-m)^2\bigr)+\lambda(x_2-x_1)$. Then
$\nabla^2F_\lambda=LI$, $\OPT=0$ is attained only at $(m,m)$, and
$F_\lambda(x)\ge\frac L2\norm{x-(m,m)}^2$ on $X$, for every $\lambda$. Let
$\gamma=\min\{y_2-y_1:y\in G_1\times G_2,\ y_2>y_1\}$. If
$\lambda>(\max d_1+\max d_2)/\gamma$, then
\[
\min\bigl\{F_\lambda(y)-d_1(y_1)-d_2(y_2):\ y\in(G_1\times G_2)\cap X\bigr\}>\OPT .
\]
\end{proposition}

\begin{proof}
On $X$, $\lambda(x_2-x_1)\ge0$, which gives the growth inequality and
$\OPT=0$ at $(m,m)$ only. A feasible grid point either has $y_1=y_2=\nu\in
G_1\cap G_2$, with corrected value $L(\nu-m)^2-d_1(\nu)-d_2(\nu)>0$ by
\eqref{eq:tu-gap}, or has $y_2-y_1\ge\gamma$, with corrected value at least
$\lambda\gamma-\max d_1-\max d_2>0$.
\end{proof}

For instance, $G_1=\{0,\frac12,1\}$, $G_2=\{0,\frac13,1\}$, $m=\frac25$ and the
corrections $d_i(\nu)=\frac L8w_i(\nu)^2$ of Section~\ref{sec:grids} satisfy
\eqref{eq:tu-gap}, since $G_1\cap G_2=\{0,1\}$,
$\frac4{25}>\frac1{32}+\frac1{72}$ and $\frac9{25}>\frac1{32}+\frac1{18}$; here
$\gamma=\frac13$, and any $\lambda>\frac{75}{288}L$ yields a corrected grid minimum
above $\OPT$. Graded grids of Section~\ref{sec:growth} built around two
different feasible centers behave in the same way: for $h=\frac1{16}$,
$\theta=\frac14$ and the centers $\frac3{10}$ and $\frac7{10}$ on $[0,1]$, one
checks that the two grids (eleven nodes each) have only $0$ and $1$ in common,
and \eqref{eq:tu-gap} holds with $m=\frac12$. The
multiplier $\lambda$ of the order constraint is invisible to the curvature
$L$ and the growth constant, and no correction depending only on them and on
the grids can absorb it. Aligned uniform grids escape because then
$G_1\cap G_2$ contains a node within $h/2$ of $m$, and
$Lh^2/4\le d_1+d_2$ there.

\begin{example}[A graded grid fails on a sum equality]\label{ex:tu-sum}
Let $X=\{x\in[0,3]^3:x_1+x_2-x_3=0\}$, which is TU with integral data, and
$F(x)=\frac L2\norm{x-(1,1,2)}^2$, so $\OPT=0$ at $(1,1,2)\in X$,
$\nabla^2F=LI$ and growth holds with $g=L/2$. Build the graded grid of
Section~\ref{sec:growth} in each coordinate around the feasible center $0$,
with $h=1$ and $\theta=\frac14$: the steps $h+\theta t$ give the common grid
$G=\{0,1,\frac94,3\}$, clipped at $3$, and the corrections
$d(0)=\frac L8$, $d(1)=d(\frac94)=\frac{25L}{128}$, $d(3)=\frac{9L}{128}$. The feasible
grid points are
$(0,0,0)$, $(1,0,1)$, $(0,1,1)$, $(\frac94,0,\frac94)$, $(0,\frac94,\frac94)$,
$(3,0,3)$ and $(0,3,3)$, since $1+1=2\notin G$ and all other sums exceed $3$ or
leave $G$. Their corrected values are
$\frac{21}8L$, $\frac{31}{64}L$, $\frac{31}{64}L$, $\frac{51}{64}L$, $\frac{51}{64}L$,
$\frac{175}{64}L$, $\frac{175}{64}L$, all positive. So the corrected grid minimum
$\frac{31}{64}L$ exceeds $\OPT$: with a sum row, a common graded pattern around a
feasible center already gives an invalid bound, without any multiplier
effect. The reason is visible in the cell $[0,1]\times[0,1]\times[1,\frac94]$:
its fiber of $x_1+x_2=x_3$ has vertices $(1,0,1)$, $(0,1,1)$ and the
non-corner vertex $(1,1,2)$, so the minimizer is not a convex combination of
feasible grid points of its cell.
\end{example}

\begin{remark}[Only coupled coordinates need uniform meshes]\label{rem:tu-hybrid}
Call a continuous coordinate \emph{coupled} if it has a nonzero coefficient in
some row of $A$, and \emph{free} otherwise. Free coordinates occur in no
row, so Lemma~\ref{lem:tu-round} can be applied to the coupled block with the
free coordinates fixed, after which the free coordinates can be rounded
independently as in Proposition~\ref{prop:cellwise}; the result is a feasible grid
point. The free coordinates may therefore use the graded grids of
Section~\ref{sec:growth}, with the coupled coordinates on the aligned
uniform mesh. This requires $\bar L$-curvature on the coupled block and
coordinate curvature on the free coordinates, and the proofs of
Lemmas~\ref{lem:inv} and~\ref{lem:states} then use only the rounding
inequality and the mesh bound of Section~\ref{sec:growth}, both of which hold.
The outcome is an XP factor $n^{p_c/2}$, where $p_c$ is the largest number of
coupled continuous coordinates in a bag, instead of $n_c^{p/2}$. We record this
only as an outline: the initial levels with mesh above $\eta$ need the same
pseudopolynomial treatment as in Theorem~\ref{thm:tu-approx}, and we have not
written out the resulting constants.
\end{remark}

\begin{remark}[Comparison with Bienstock and Mu\~noz]\label{rem:tu-bm}
Bienstock and Mu\~noz \cite[Theorems~4 and~15]{BienstockMunoz2018} construct,
for mixed-binary polynomial problems with variables scaled to $[0,1]$, a linear
objective $c$, constraints of degree at most $\pi$ and coefficient $1$-norm at
most $\mathfrak F$, and a constraint intersection graph of treewidth $\omega$, a
linear program with $O\bigl((2\pi/\epsilon)^{\omega+1}n\log(\pi/\epsilon)\bigr)$
variables and constraints whose solutions violate each constraint by at most
$\mathfrak F\epsilon$ and are $\lVert c\rVert_1\epsilon$-optimal. They also argue that,
unless $\mathrm P=\mathrm{NP}$, neither the pseudopolynomial dependence on
$1/\epsilon$ nor the coefficient scaling of the violation can be removed for that
class. The results here differ in four respects. (i) Feasibility is exact, but
only for TU continuous columns with aligned data; their constraints are
arbitrary polynomials. (ii) A nonlinear objective is handled directly through
$\bar L$, whereas their objective is linear and nonlinear objectives are
encoded by additional constraints, which may increase the width. (iii) Under
growth with finitely many optimal coordinate values, the per-level tables do
not depend on $\epsilon$, and $\epsilon$ enters only through $O(\log(1/\epsilon))$ levels;
this does not conflict with their lower bounds, which concern their whole
class and hence instances without growth or TU structure, and the factor
$(n_c\kappa_c)^{p/2}$ can itself be exponential in $I$ because $\kappa_c$ can be. (iv) Without growth the last level of
TU-GRID has $O\bigl((W\sqrt{n_c\bar L/\epsilon}+1)^p\bigr)$ table entries, an exponent
$p/2$ in $1/\epsilon$ against their $\omega+1$ for comparable decompositions; the
halving comes from second-order rounding error with exact feasibility, and
the comparison is only indicative because the tolerances are normalized
differently. Neither result contains the other.

Proximity scaling for separable convex objectives over constraint matrices
with bounded subdeterminants \cite{HochbaumShanthikumar1990} also refines
aligned grids inside windows around the previous scaled solution, with a
window radius proportional to $n$ times the scale for TU matrices. There the
subproblems are convex and solved by linear programming. Here the objective is
nonseparable and nonconvex, the subproblems are solved by dynamic programming
over the decomposition, and the window radius $\sqrt{n_c\kappa_c}\,h_j$ comes
from growth and the rounding allowance.
\end{remark}
